\documentclass[lettersize,journal]{IEEEtran}
\usepackage{amsmath,amsfonts}
\usepackage{algorithmic}
\usepackage{algorithm}
\usepackage{array}
\usepackage{multirow}
\usepackage[caption=false,font=normalsize,labelfont=sf,textfont=sf]{subfig}
\usepackage{textcomp}
\usepackage{stfloats}
\usepackage{url}
\usepackage{verbatim}
\usepackage{graphicx}
\usepackage{cite}
\usepackage{lipsum}
\hyphenation{op-tical net-works semi-conduc-tor IEEE-Xplore}
% updated with editorial comments 8/9/2021

\begin{document}

\title{LE-MVCL: Label-Efficient Encrypted Traffic Classification via Asymmetric Multi-View Contrastive Learning}

\author{Hanif Nur Ilham Sanjaya, Alfan Presekal}

% The paper headers
\markboth{IEE Transaction Journal \LaTeX\ Class Files,~Vol.~xx, No.~xx, January~2026}%
{Shell \MakeLowercase{\textit{et al.}}: A Sample Article Using IEEEtran.cls for IEEE Journals}

\IEEEpubid{0000--0000/00\$00.00~\copyright~2021 IEEE}
% Remember, if you use this you must call \IEEEpubidadjcol in the second
% column for its text to clear the IEEEpubid mark.

\maketitle

\begin{abstract}
The rapid proliferation of encrypted tunneling protocols, particularly Virtual Private Networks (VPNs), has rendered traditional Deep Packet Inspection (DPI) ineffective, creating an urgent need for classification methods that are both robust and label-efficient. To address this challenge, we propose \textbf{LE-MVCL} (Label-Efficient Encrypted Traffic Classification via Asymmetric Multi-View Contrastive Learning), a novel semi-supervised framework. Unlike standard contrastive approaches that rely on symmetric data augmentations, we introduce an \textit{Asymmetric Multi-View Pre-training} strategy where a Deep MLP processing global flow statistics acts as a "semantic guide" for a 1D-CNN encoder learning from raw payload bytes. This aligns the semantic representation of the payload with robust behavioral signatures without requiring manual labels. In the fine-tuning stage, we employ a \textit{Hybrid Late-Fusion} architecture, concatenating learned semantic embeddings with a comprehensive set of raw feature views—specifically \textit{Alpha} (Payload), \textit{Beta} (Flow Dynamics), \textit{Gamma} (Burst Behavior), and \textit{FFT} (Spectral Analysis)—to explicitly preserve granular domain properties often lost in pure neural encodings. Experimental results demonstrate superior label efficiency, with the model achieving a binary detection F1-score of \textbf{92.2\% using only 10\% labeled data} (scaling to 94.9\% with 30\% supervision), while effectively performing hierarchical classification of specific applications with \textbf{$>$91\% F1}, significantly mitigating the label scarcity problem in encrypted traffic analysis.
\end{abstract}

\begin{IEEEkeywords}
Contrastive Learning, Self-Supervised Learning, Label Efficiency, Encrypted Traffic Classification, VPN Detection.
\end{IEEEkeywords}

\section{Introduction}
\IEEEPARstart{T}{he} rapid evolution of internet privacy mechanisms has led to the ubiquitous adoption of encryption protocols. As noted by Kahn \cite{kahn1996codebreakers}, encryption techniques have advanced significantly to meet the demand for secure communication across military, commercial, and personal domains. Among these technologies, Virtual Private Networks (VPNs) have become a fundamental tool for users seeking to bypass geo-restrictions, secure remote access, and safeguard sensitive data from unauthorized surveillance \cite{singhpall1996vpn}. However, while VPNs provide essential privacy benefits, they also introduce significant challenges for network management and security monitoring. The anonymity provided by VPNs can be exploited for malicious activities, creating a complex "dual-use" dilemma where legitimate privacy preservation conflicts with the need to detect cyber threats \cite{dipietro2021new}.

Traditionally, network traffic classification relied on port-based methods, which matched packet headers to the Internet Assigned Numbers Authority (IANA) registry \cite{yoon2009internet}. However, empirical studies have shown that port-based classification is increasingly unreliable, with failure rates ranging from 30\% to 70\% due to the use of dynamic port negotiation and obfuscation techniques \cite{moore2005toward}. To address this, Deep Packet Inspection (DPI) emerged as a robust alternative, analyzing packet payloads for specific signatures \cite{finsterbusch2013survey}. Yet, the widespread deployment of modern encryption standards renders DPI largely ineffective, as payload content is obfuscated and unreadable without decryption keys \cite{wang2017endtoend}.

Consequently, the research community has shifted toward Machine Learning (ML) and Deep Learning (DL) paradigms. Recent surveys indicate a clear trend moving from manual feature engineering toward DL models that utilize raw packet data or time-series representations \cite{razooqi2025vpn}. Approaches such as "Deep Packet" \cite{lotfollahi2019deep} and 1D-CNNs \cite{wang2017endtoend} have demonstrated state-of-the-art performance by automatically learning complex patterns from encrypted flows. However, these supervised DL models face two critical limitations. First, they require massive amounts of labeled data, which is labor-intensive and expensive to acquire in real-world environments \cite{menezes2018handbook}. Second, existing semi-supervised approaches \cite{lin2022semisupervised} often fail to effectively align the multi-modal nature of network traffic. They typically treat payload signatures and statistical flow behaviors as separate entities, missing the opportunity to use one view to semantically guide the representation learning of the other.

To bridge this gap, we propose LE-MVCL (Label-Efficient Multi-View Contrastive Learning), a novel framework designed to achieve robust encrypted traffic classification under severe label scarcity. LE-MVCL distinguishes itself through an \textit{asymmetric} pre-training strategy that leverages unlabeled data to align a Deep 1D-CNN (Payload View) with a Deep MLP (Statistical View). Furthermore, we introduce a comprehensive feature synthesis encompassing payload, flow dynamics, and burst behaviors, ensuring that granular statistical precision is preserved alongside deep semantic embeddings.

The main contributions of this paper are as follows:
\begin{itemize}
    \item The LE-MVCL Framework: We propose a novel framework that adapts Contrastive Learning to the domain of VPN identification. Unlike standard supervised methods, our approach leverages unlabeled data to learn robust flow representations, achieving 94.5\% Binary F1 with only 10\% labeled data and 96.0\% with 30\% labeled data, significantly mitigating the label scarcity problem.
    
    \item Asymmetric Multi-View Pre-training: We introduce an asymmetric contrastive pre-training strategy that aligns a Deep 1D-CNN (Payload View) with a Deep MLP (Statistical View). We demonstrate that using the statistical view as a "semantic guide" during pre-training forces the CNN to learn high-level behavioral patterns from raw bytes, significantly outperforming standard end-to-end training.
    
    \item Hybrid Late-Fusion Architecture: We empirically demonstrate that a Hybrid Late-Fusion strategy—combining learned semantic embeddings with raw statistical features—yields superior performance compared to pure deep learning approaches. This validates that preserving the granular precision of statistical metrics is critical for distinguishing fine-grained VPN applications (e.g., Skype vs. Hangouts).
    
    \item Robust Feature Engineering: We synthesize a comprehensive multi-view feature set comprising the Alpha (Payload), Beta (Flow Dynamics), Gamma (Burst Behavior), and FFT (Spectral) components. We prove that fusing these diverse views enables the detection of obfuscated traffic patterns, maintaining high efficacy ($>$93\% App F1) even when distinguishing between isotocol services.
\end{itemize}
\section{Related Work}
\label{sec:related_work}



Encrypted traffic classification has evolved substantially over the past decade, transitioning from traditional port-based heuristics toward data-driven deep learning approaches. The widespread adoption of encryption protocols and tunneling mechanisms, particularly Virtual Private Networks (VPNs), has rendered signature-based inspection and port-based identification increasingly ineffective. Consequently, modern research focuses on learning discriminative representations directly from traffic data. Nevertheless, two fundamental challenges remain: the scarcity of labeled datasets and the difficulty of effectively integrating heterogeneous traffic representations such as raw payloads and flow-level statistics. Table \ref{tab:rw_comparison} summarizes the key differences between existing frameworks and the proposed LE-MVCL approach.

Early deep learning methods demonstrated that raw packet payloads could be treated as sequential data and processed using neural architectures originally developed for other domains. For instance, Wang \textit{et al.} \cite{wang2017end} and Lotfollahi \textit{et al.} \cite{lotfollahi2020deep} showed that 1D-Convolutional Neural Networks (1D-CNNs) and stacked autoencoders can automatically learn hierarchical representations from raw byte sequences, effectively modeling network traffic in a manner similar to natural language sequences. Subsequent work extended this paradigm by incorporating temporal modeling through Long Short-Term Memory (LSTM) networks \cite{yao2019identification} and structural modeling through Graph Neural Networks (GNNs) \cite{huoh2022flow}. These approaches achieved strong performance on benchmark datasets such as ISCX-VPN-2016. However, their success heavily depends on large-scale labeled datasets, which are costly and difficult to obtain in real-world operational environments such as Security Operations Centers (SOCs). In practice, obtaining ground-truth labels for encrypted tunnels often requires manual inspection and controlled experiments, resulting in limited labeled data. As a consequence, purely supervised approaches typically suffer from performance degradation when deployed in low-resource or dynamically evolving network environments.

To alleviate the dependence on labeled data, Self-Supervised Learning (SSL) has recently attracted increasing attention in encrypted traffic analysis. Inspired by representation learning techniques developed in computer vision, frameworks such as SimCLR \cite{chen2020simple} have been adapted to network traffic classification. These approaches typically employ symmetric contrastive learning, where two augmented views of the same traffic sample are generated through stochastic transformations such as packet masking, byte shuffling, or cropping. A shared encoder is then trained to maximize the agreement between these augmented representations. While effective in image domains, such symmetric augmentation strategies are less suitable for encrypted network traffic. Unlike visual data, encrypted packets follow strict protocol structures, and arbitrary byte-level perturbations may violate the syntactic integrity of encapsulation protocols. 

This limitation is illustrated in Fig. \ref{fig:vpn_syntax}, which visualizes the impact of symmetric masking on real OpenVPN packets from the ISCXVPN2016 dataset. In a valid OpenVPN packet, the outer header contains critical fields such as the OpCode and Sequence Identifier that ensure tunnel consistency and anti-replay protection. Random byte masking may corrupt these fields, producing packet representations that do not correspond to any legitimate protocol state. As a result, the encoder may be forced to learn representations from logically inconsistent traffic patterns, which can negatively affect downstream classification performance. In contrast, LE-MVCL preserves the original byte sequence and derives the contrastive learning signal from complementary views of the same flow.

\begin{table*}[!htbp]
\caption{Comparison of LE-MVCL with State-of-the-Art Frameworks}
\label{tab:rw_comparison}
\centering
\begin{tabular}{l c c c c}
\hline
\textbf{Framework} & \textbf{Learning Paradigm} & \textbf{Pre-training Strategy} & \textbf{Fusion Type} & \textbf{Label Efficiency} \\
\hline
Deep Packet \cite{lotfollahi2020deep} & Supervised & None & None (Payload Only) & Low \\
DFR \cite{zeng2019deep} & Supervised & None & Deep / Latent & Low \\
Flow-GNN \cite{huoh2022flow} & Supervised & None & Graph + Flow & Medium \\
Generic SSL (SimCLR-Net) & Self-Supervised & Symmetric (Augmentations) & None & Medium \\
\textbf{LE-MVCL (Ours)} & \textbf{Semi-Supervised} & \textbf{Asymmetric (Payload $\leftrightarrow$ Stats)} & \textbf{Hybrid Late-Fusion} & \textbf{High} \\
\hline
\end{tabular}
\end{table*}

\begin{figure}[h]
    \centering
    \includegraphics[width=\linewidth]{img/vpn_syntax_destruction.png}
    \caption{Impact of Symmetric Masking on Encrypted VPN Traffic. 
    Visualized using a real OpenVPN UDP packet from the ISCXVPN2016 dataset. 
    (A) LE-MVCL Input: The raw byte sequence is preserved, maintaining the validity of the outer VPN header (OpCode \texttt{0x1D}) and Sequence ID. 
    (B) Symmetric Masking: Random byte masking (standard in SimCLR) corrupts critical tunnel fields. Masking the \textit{OpCode} renders the packet protocol-invalid, while masking the \textit{Sequence ID} creates artifacts resembling replay attacks. This forces the encoder to learn from logically broken data.}
    \label{fig:vpn_syntax}
\end{figure}

Another limitation of existing SSL-based traffic analysis methods is their inability to exploit the inherent multi-view nature of network flows. Encrypted traffic simultaneously exhibits low-level payload structures and high-level statistical patterns, such as burst behavior, packet inter-arrival dynamics, and spectral characteristics. Many existing frameworks focus primarily on a single representation or treat multiple modalities through symmetric augmentation, thereby underutilizing the complementary information available across different traffic views.

To address these limitations, LE-MVCL introduces an \textit{Asymmetric Multi-View Pre-training} strategy that explicitly models the dual nature of encrypted traffic. Instead of generating artificial augmentations from a single modality, the framework aligns two complementary representations of the same flow. A 1D-CNN encoder processes the raw payload bytes, while a multilayer perceptron (MLP) processes global flow-level statistical features. The statistical representation acts as a semantic guide, encouraging the payload encoder to learn embeddings that correlate with stable behavioral patterns of the traffic flow. This asymmetric design enables the model to capture higher-level protocol semantics without requiring manual labels or syntactically destructive augmentations.

In addition to representation learning, another critical design consideration in encrypted traffic classification is the fusion of heterogeneous features. Several recent studies have explored multimodal architectures that combine payload-based features with statistical descriptors. Frameworks such as Deep-Full-Range (DFR) \cite{zeng2019deep} and related approaches \cite{liu2024adaptive} typically employ deep fusion mechanisms, where multiple feature modalities are first projected into a shared latent space through neural encoders before classification. Although this strategy allows joint representation learning, it may inadvertently compress fine-grained statistical signals. Metrics such as precise packet inter-arrival intervals, burst patterns, or spectral characteristics can become smoothed or distorted when passed through multiple nonlinear transformations, potentially reducing their discriminative value for distinguishing closely related applications.

To mitigate this issue, LE-MVCL adopts a \textit{Hybrid Late-Fusion} architecture. Instead of projecting all modalities into a shared deep latent space, the semantic embeddings learned from the payload encoder are concatenated directly with raw normalized statistical features, including Alpha, Beta, Gamma, and FFT-based spectral descriptors. This design preserves the high-resolution statistical information while still benefiting from the semantic abstraction learned through deep representation learning. As a result, the final classifier can leverage both the expressive capability of deep embeddings and the precise discriminative power of handcrafted statistical features, leading to improved performance in fine-grained encrypted application classification.





% \section{Related Work}
% \label{sec:related_work}

% \begin{table*}[htbp]
% \caption{Comparison of LE-MVCL with State-of-the-Art Frameworks}
% \label{tab:rw_comparison}
% \centering
% \begin{tabular}{l c c c c}
% \hline
% \textbf{Framework} & \textbf{Learning Paradigm} & \textbf{Pre-training Strategy} & \textbf{Fusion Type} & \textbf{Label Efficiency} \\
% \hline
% Deep Packet \cite{lotfollahi2020deep} & Supervised & None & None (Payload Only) & Low \\
% DFR \cite{zeng2019deep} & Supervised & None & Deep / Latent & Low \\
% Flow-GNN \cite{huoh2022flow} & Supervised & None & Graph + Flow & Medium \\
% Generic SSL (SimCLR-Net) & Self-Supervised & Symmetric (Augmentations) & None & Medium \\
% \textbf{LE-MVCL (Ours)} & \textbf{Semi-Supervised} & \textbf{Asymmetric (Payload $\leftrightarrow$ Stats)} & \textbf{Hybrid Late-Fusion} & \textbf{High} \\
% \hline
% \end{tabular}
% \end{table*}

% The classification of encrypted traffic has undergone a significant paradigm shift, evolving from port-based heuristics to sophisticated Deep Learning (DL) architectures. However, the field currently faces dual challenges: the prohibitive cost of acquiring labeled datasets and the technical difficulty of effectively fusing heterogeneous traffic views. In this section, we situate the proposed LE-MVCL framework within the broader landscape of supervised deep learning, contrastive representation learning, and multimodal fusion strategies.

% \subsection{Supervised Deep Learning and the Label Scarcity Bottleneck}
% With the obsolescence of port-based classification and the decreasing efficacy of traditional statistical machine learning, Deep Learning has emerged as the dominant paradigm for encrypted traffic analysis. Pioneering works by Wang et al. \cite{wang2017end} and Lotfollahi et al. \cite{lotfollahi2020deep} demonstrated that 1D-Convolutional Neural Networks (1D-CNNs) and Stacked Autoencoders (SAE) could automatically learn feature hierarchies from raw packet bytes, effectively treating network traffic analysis as a sequence modeling task akin to Natural Language Processing. Subsequent research has expanded this capability by incorporating Long Short-Term Memory (LSTM) networks \cite{yao2019identification} to capture temporal dependencies and Graph Neural Networks (GNN) \cite{huoh2022flow} to model the topological structure of traffic flows.

% Despite their high accuracy on benchmark datasets like ISCX-VPN-2016, these supervised models are fundamentally constrained by their reliance on massive, fully labeled datasets. In realistic operational environments, such as Security Operations Centers (SOCs), acquiring ground-truth labels for encrypted tunnels is a labor-intensive and costly process. Consequently, standard supervised architectures often suffer from severe performance degradation in low-resource settings. This limitation underscores the critical need for label-efficient frameworks capable of learning robust representations from the vast amounts of unlabeled traffic available in production networks.

% \subsection{From Symmetric to Asymmetric Contrastive Learning}
% To address the label scarcity problem, Self-Supervised Learning (SSL) has gained traction within the networking community. Techniques originally developed for Computer Vision, such as SimCLR \cite{chen2020simple}, have been adapted to traffic analysis, typically employing a \textit{symmetric} contrastive learning approach. In these methods, a single encoder is trained to maximize the agreement between two augmented views of the same flow, generated via techniques like random packet masking or payload shuffling.

% We argue, however, that symmetric augmentation is suboptimal for encrypted tunnel traffic. Unlike images, where pixel noise is tolerable, randomly masking packet bytes can destroy the rigid syntax of encapsulation protocols (e.g., OpenVPN, IPsec). In a VPN context, the "payload" seen by the model begins with a delicate outer header containing Opcodes and Sequence IDs required for tunnel integrity.

% To illustrate this empirically, we visualize the impact of these augmentations on real OpenVPN traffic from the ISCXVPN2016 dataset in Fig. \ref{fig:vpn_syntax}.

% \begin{figure}[h]
%     \centering
%     \includegraphics[width=\linewidth]{img/vpn_syntax_destruction.png}
%     \caption{Impact of Symmetric Masking on Encrypted VPN Traffic. 
%     Visualized using a real OpenVPN UDP packet from the ISCXVPN2016 dataset. 
%     (A) LE-MVCL Input: The raw byte sequence is preserved, maintaining the validity of the outer VPN header (OpCode \texttt{0x1D}) and Sequence ID. 
%     (B) Symmetric Masking: Random byte masking (standard in SimCLR) corrupts critical tunnel fields. Masking the \textit{OpCode} renders the packet protocol-invalid, while masking the \textit{Sequence ID} creates artifacts resembling replay attacks. This forces the encoder to learn from logically broken data.}
%     \label{fig:vpn_syntax}
% \end{figure}

% As demonstrated in Fig. \ref{fig:vpn_syntax}, a standard stochastic augmentation that blindly zeros out bytes has a high probability of corrupting the \textit{VPN OpCode} (identifying the packet type) or the \textit{Sequence Number} (critical for anti-replay logic). Unlike image cropping, which leaves the semantic object intact, this transformation renders the packet syntactically invalid—effectively asking the model to learn a representation for a "broken" tunnel state that never occurs in legitimate traffic. LE-MVCL (Fig. \ref{fig:vpn_syntax}A) addresses this by preserving the exact byte sequence and instead deriving the contrastive signal from the flow's global behavioral statistics.

% Furthermore, existing network SSL methods often ignore the inherent duality of traffic, failing to exploit the stable statistical properties of a flow to guide the learning of the more volatile payload representation. 

% In contrast, LE-MVCL introduces an \textit{Asymmetric Multi-View Pre-training} strategy. rather than relying on artificial augmentations, we align two distinct deep encoders: a 1D-CNN processing the raw payload and an MLP processing global flow statistics. By treating the robust statistical view as a "Semantic Guide," our approach forces the CNN to learn representations that correlate with global behavioral patterns, thereby capturing high-level protocol semantics without requiring explicit supervision.

% \subsection{Deep Embedding vs. Hybrid Late-Fusion}
% Recognizing that neither payload nor statistics alone suffice for fine-grained classification, recent literature has moved towards hybrid architectures. Frameworks such as "Deep-Full-Range" (DFR) \cite{zeng2019deep} and other multimodal approaches \cite{liu2024adaptive} typically employ a "Deep Fusion" strategy. In this paradigm, statistical features are passed through deep neural encoders and projected into a shared latent space alongside payload embeddings before classification.

% While conceptually appealing, we empirically identify a critical flaw in Deep Fusion regarding \textit{statistical compression}. Projecting precise metrics—such as exact inter-arrival times or spectral entropy peaks—through deep non-linear layers often "smooths out" the granular anomalies necessary to distinguish between isotocol services (e.g., distinguishing a Skype voice call from a Hangouts voice call). 

% To overcome this, LE-MVCL adopts a \textit{Hybrid Late-Fusion} architecture. Instead of fusing abstract latent vectors, we concatenate the learned semantic embeddings from the CNN directly with the \textit{raw, normalized} statistical features (Alpha, Beta, Gamma, FFT). This strategy ensures that the granular precision of the statistical view is preserved intact for the final decision boundary, enabling superior performance in fine-grained application classification where traditional deep fusion models often fail.
\section{Methodology}
\label{sec:methodology}

\subsection{Dataset Curation and Preprocessing}
\label{subsec:dataset_curation}

To ensure robust generalization across different network environments, we utilized two distinct encrypted traffic datasets: the widely-used ISCX-VPN-2016 benchmark and the VNAT dataset.

\subsubsection{Data Sources}
\begin{itemize}
    \item ISCX-VPN-2016 \cite{draper2016characterization}: The primary benchmark for encrypted traffic analysis, containing real-world traffic generated from Alice and Bob scenarios. It encompasses 7 specific traffic categories (e.g., Browsing, Email, Streaming) and 16 granular applications captured over both standard and VPN channels.
    \item VNAT (VPN/Non-VPN Application Traffic) \cite{vnat_dataset}: To augment the diversity of our training corpus, we incorporated the VNAT dataset from MIT Lincoln Laboratory. This dataset includes PCAP captures of both encrypted and unencrypted flows across 10 distinct applications (including SSH, RDP, and Zoiper) generated in a controlled virtual subnet environment, providing high-fidelity samples for background traffic and specific tunneling behaviors.
\end{itemize}

% Use figure* (with an asterisk) for full-width images
\begin{figure*}[ht] 
    \centering
    % width=\textwidth fills the printable area from left margin to right margin
    \includegraphics[width=\textwidth]{iscx_dist.png}
    \caption{Distribution of the clean ISCX dataset.}
    \label{fig:iscx_dist}
\end{figure*}

% Use figure* (with an asterisk) for full-width images
\begin{figure*}[ht] 
    \centering
    % width=\textwidth fills the printable area from left margin to right margin
    \includegraphics[width=\textwidth]{vnat_dist.png}
    \caption{Distribution of the clean VNAT dataset.}
    \label{fig:vnat_dist}
\end{figure*}

\begin{table}[htbp]
\caption{Summary of Preprocessing Rules}
\label{tab:preprocessing_rules}
\centering
\begin{tabular}{l l l}
\hline
\textbf{Filter Type} & \textbf{Condition} & \textbf{Action} \\
\hline
Protocol & DNS, NTP, DHCP & Drop \\
Port Specific & UDP 137 (NBNS), UDP 5355 (LLMNR) & Drop \\
TCP Integrity & First Packet $\neq$ SYN & Drop \\
UDP Noise & DstIP=255.255.255.255 AND Payload $\in$ 'Beacon$\sim$' & Drop \\
Min Length & Packet Count $< 2$ & Drop \\
\hline
\end{tabular}
\end{table}

\subsubsection{Preprocessing Pipeline}
Raw packet captures were first segmented into bidirectional flows based on the standard 5-tuple identifier (Source IP, Destination IP, Source Port, Destination Port, Protocol) using the SplitCap tool. To ensure data quality and comparability with state-of-the-art benchmarks, we strictly implemented the cleaning and filtering protocol defined by Park et al. \cite{park2024fast}. The pipeline consists of four critical stages that are summarized in Table \ref{tab:preprocessing_rules}.

\begin{enumerate}
    \item Noise \& Background Traffic Removal:
    We filtered out infrastructure traffic irrelevant to the target application classification. This includes removing protocols such as DNS (Port 53), NTP (Port 123), and DHCP (Port 67/68). Additionally, consistent with \cite{park2024fast}, we explicitly dropped NetBIOS Name Service (NBNS) traffic on UDP port 137 and Link-Local Multicast Name Resolution (LLMNR) on UDP port 5355.

    \item TCP Connection Sanity Check:
    To ensure the integrity of TCP flows used for payload analysis, we discarded any TCP flow that did not begin with a valid handshake initiation. Specifically, we filtered out flows where the first packet did not contain the \texttt{SYN} flag (Flag 'S'), eliminating incomplete or mid-stream captured connections.

    \item Beacon \& Broadcast Filtering:
    We observed specific recurrent UDP broadcast patterns that do not carry user-generated content. We applied a rule to remove UDP flows directed to the broadcast IP \texttt{255.255.255.255} that contained the specific payload signature ``\texttt{Beacon$\sim$}''. This step is crucial for removing local discovery noise that can skew statistical feature distributions.
    
    \item Short Flow Pruning:
    Flows containing fewer than 2 packets were discarded as they lack sufficient information for both the temporal statistical view (Inter-Arrival Times) and the payload contrastive view.
\end{enumerate}

Following this rigorous filtration process, we obtained a final dataset that is based on its distribution as shown in Fig. \ref{fig:iscx_dist} and \ref{fig:vnat_dist}.


% ==========================

\subsection{Multi-View Feature Generation}
\label{subsec:feature_generation}

To overcome the limitations of single-view classification, the LE-MVCL framework constructs a composite representation of network traffic by fusing learned semantic representations with engineered statistical fingerprints, following the common multimodal design of jointly leveraging payload-derived and flow-level information for encrypted traffic classification \cite{zeng2019deep,chen2025mff}. Let a network flow $\mathcal{F}$ be defined as a sequence of packets $\mathcal{F} = \{p_1, p_2, \dots, p_K\}$, where each packet $p_i$ is a tuple $(h_i, \mathbf{b}_i, t_i)$ consisting of header information $h_i$, a payload byte sequence $\mathbf{b}_i$, and an arrival timestamp $t_i$. This multi-view approach addresses the duality of encrypted traffic, where the payload carries protocol-specific signatures ($\alpha$) while the flow dynamics reveal application-layer behaviors ($\beta, \gamma, \beta'_{FFT}$).

\subsubsection{View 1: Learned Payload Representation (The Alpha Component)}
\label{subsubsec:alpha_view}

The first view, termed the ``Alpha'' component ($\alpha$), functions as a semantic extractor designed to identify non-linear dependencies in the protocol handshake—such as TLS Client Hello extensions or cipher suite negotiations—that remain invisible to pure time-series analysis. This viewpoint is consistent with recent datagram-level representation learning approaches that directly model encrypted payload bytes with deep encoders (e.g., pre-trained transformers) \cite{lin2022etbert}. This viewpoint is consistent with recent datagram-level representation learning approaches that directly model encrypted payload bytes with deep encoders (e.g., pre-trained transformers) \cite{lin2022etbert}.

Following the standard Deep Packet Inspection (DPI) methodology established in \cite{lotfollahi2020deep}, we isolate the first $N=10$ packets of the flow. Let $\mathbf{B}_{raw}$ denote the concatenated sequence of payload bytes from these packets:
\begin{equation}
    \mathbf{B}_{raw} = \bigoplus_{i=1}^{N} \mathbf{b}_i
\end{equation}
where $\oplus$ denotes the concatenation operator. To ensure a uniform input dimension for the neural network, $\mathbf{B}_{raw}$ is strictly truncated or zero-padded to a fixed length $L=784$ bytes. This specific window size is selected to capture the critical initial handshake phase where unencrypted metadata, such as the Server Name Indication (SNI), is typically exchanged \cite{wang2017end} \cite{rfc6066}.

The raw integer byte sequence $\mathbf{x} \in \{0, \dots, 255\}^L$ is normalized to the interval $[0,1]$ via min-max scaling:
\begin{equation}
    \hat{x}_j = \frac{x_j}{255}, \quad \forall j \in \{1, \dots, L\}
\end{equation}

Rather than using these raw bytes directly, we pass the normalized sequence $\hat{\mathbf{x}}$ through a pre-trained 1D-Convolutional Neural Network (1D-CNN) encoder, denoted as $f_{\theta}(\cdot)$. This encoder compresses the high-dimensional input into a dense latent vector $\mathbf{z}_{\alpha} \in \mathbb{R}^{128}$:
\begin{equation}
    \mathbf{z}_{\alpha} = f_{\theta}(\hat{\mathbf{x}})
\end{equation}
This vector $\mathbf{z}_{\alpha}$ effectively distills the sparse raw bytes into a compact semantic signature associated with the specific application protocol.

\subsubsection{View 2: Global Statistical Signatures (The Greek Components)}
\label{subsubsec:greek_view}

While encryption effectively masks payload content, it cannot fully obfuscate global traffic patterns. To capture these leakage channels, we engineer a set of statistical features collectively termed the ``Greek'' components. This perspective is aligned with recent flow-level sequence modeling work that explicitly targets packet-to-packet dynamics within a flow \cite{chen2024miett}. This perspective is aligned with recent flow-level sequence modeling work that explicitly targets packet-to-packet dynamics within a flow \cite{chen2024miett}.

\paragraph{Flow Dynamics ($\beta$)}
The foundation of the statistical view is the $\beta$ component. Let $S = \{|p_1|, \dots, |p_K|\}$ be the sequence of packet sizes and $T = \{\Delta t_1, \dots, \Delta t_{K-1}\}$ be the sequence of Inter-Arrival Times (IAT), where $\Delta t_i = t_{i+1} - t_i$. We define a statistical extraction function $\Phi(\cdot)$ that computes the first four statistical moments (Mean $\mu$, Variance $\sigma^2$, Skewness $\gamma_1$, Kurtosis $\kappa$) along with the Min and Max values:
\begin{equation}
    \Phi(\mathbf{v}) = \left[ \min(\mathbf{v}), \max(\mathbf{v}), \mu_{\mathbf{v}}, \sigma_{\mathbf{v}}, \frac{\mathbb{E}[(\mathbf{v}-\mu)^3]}{\sigma^3}, \frac{\mathbb{E}[(\mathbf{v}-\mu)^4]}{\sigma^4} \right]
\end{equation}
The final feature vector $\mathbf{v}_{\beta}$ is the concatenation of these statistics computed bi-directionally (forward and backward flows) for both $S$ and $T$, distinguishing the request-response nature of interactive protocols \cite{draper2016characterization}.

\paragraph{Burst Behavior ($\gamma$)}
To model application-layer buffering strategies, we employ the $\gamma$ component. We define a ``burst'' as a subsequence of packets arriving with high frequency. A set of packets $\{p_i, \dots, p_{i+m}\}$ forms a burst $B_j$ if the IAT between all adjacent packets in the sequence is below a dynamic threshold $\tau$ (typically $\tau < 1.0s$):
\begin{equation}
    B_j = \{ p_k \mid \Delta t_{k-1} < \tau \}
\end{equation}
From the set of all bursts $\mathcal{B} = \{B_1, \dots, B_M\}$, we extract features such as Burst Density (bytes per burst) and the Active-to-Idle Ratio ($R_{AI}$):
\begin{equation}
    R_{AI} = \frac{\sum_{j=1}^{M} \text{duration}(B_j)}{\text{TotalFlowDuration} - \sum_{j=1}^{M} \text{duration}(B_j)}
\end{equation}
This view specifically targets the "ON-OFF" transmission patterns characteristic of video buffering \cite{dubin2017experimental}.

\paragraph{Spectral Analysis ($\beta'_{FFT}$)}
Finally, to detect automated behavior hidden within the noise, we utilize the $\beta'_{FFT}$ component. We apply the Discrete Fourier Transform (DFT) to the sequence of Packet Sizes $S$ to convert the time-domain signal into the frequency domain $X[k]$:
\begin{equation}
    X[k] = \sum_{n=0}^{N-1} s_n \cdot e^{-i 2\pi k n / N}
\end{equation}
From the resulting power spectrum $|X[k]|^2$, we extract the Top-$K$ dominant frequencies and the Spectral Entropy ($H_{spec}$), defined as:
\begin{equation}
    H_{spec} = -\sum_{k=1}^{K} P(k) \log_2 P(k), \quad \text{where } P(k) = \frac{|X[k]|^2}{\sum_j |X[j]|^2}
\end{equation}
This frequency-domain analysis identifies periodic "heartbeat" signals often generated by VPN keep-alive packets, appearing as distinct peaks (low entropy) in the frequency spectrum \cite{shin2020encrypted}.

\subsection{Phase I: Asymmetric Multi-View Pre-training}
\label{subsec:phase1_contrastive}

To surmount the limitations of label scarcity, LE-MVCL introduces an asymmetric self-supervised pre-training strategy. We utilize a dual-branch architecture designed to align the latent representations of the \textit{Alpha} component (raw payload) with the diverse statistical properties of the \textit{Beta}, \textit{Gamma}, and \textit{FFT} components. The core objective is not merely alignment, but using the high-level statistical view as a "semantic guide" to force the CNN to learn robust behavioral patterns from raw bytes. The complete pre-training procedure is outlined in Algorithm \ref{alg:pretraining}.

\subsubsection{Architecture Overview}
The framework consists of two parallel encoders with distinct inductive biases:

\begin{itemize}
    \item Branch A (Payload View - Alpha): A Deep 1D-CNN processes the raw payload sequence $x_{\alpha} \in \mathbb{R}^{1024}$ to produce a semantic representation $h_{cnn}$. This is followed by a non-linear projection head (Dense $\rightarrow$ ReLU $\rightarrow$ Dense) to map the representation to a lower-dimensional embedding $z_{pay} \in \mathbb{R}^{64}$.
    
    \item Branch B (Statistical View - Beta/Gamma/FFT): A Deep MLP encoder processes the concatenated statistical components $x_{stat} = [x_{\beta}, x_{\gamma}, x_{fft}] \in \mathbb{R}^{N_{stat}}$. It comprises a dense layer (256 units, ReLU), a Dropout layer ($p=0.3$) for regularization, and a second dense layer (128 units). A symmetric projection head maps this to $z_{stat} \in \mathbb{R}^{64}$.
\end{itemize}

\subsubsection{Objective Function (Asymmetric NT-Xent)}
We employ the Normalized Temperature-scaled Cross Entropy (NT-Xent) loss to maximize the mutual information between the payload and statistical views. For a minibatch of $N$ flows, we treat the statistical view $z_{stat}$ as the anchor to guide the payload view $z_{pay}$. For a given flow $i$, the loss function is defined as:

\begin{equation}
    \ell_{i,j} = -\log \frac{\exp(\text{sim}(z_{pay}^{(i)}, z_{stat}^{(i)}) / \tau)}{\sum_{k=1}^{2N} \mathbb{1}_{[k \neq i]} \exp(\text{sim}(z_{pay}^{(i)}, z_{k}) / \tau)}
\end{equation}

where $\text{sim}(u, v) = u^T v / \|u\| \|v\|$ denotes cosine similarity, $\mathbb{1}$ is the indicator function, and $\tau=0.1$ is the temperature parameter. By minimizing this loss, we enforce the 1D-CNN to implicitly learn the "global" flow dynamics (e.g., burstiness encoded in Branch B) solely from the "local" packet bytes, effectively transferring knowledge from the statistical domain to the payload domain without manual labels.

\begin{algorithm}[H]
\caption{LE-MVCL Asymmetric Pre-training Strategy}
\label{alg:pretraining}
\begin{algorithmic}[1]
\REQUIRE Unlabeled Dataset $\mathcal{D} = \{ (x_{\alpha}^{(i)}, x_{stat}^{(i)}) \}_{i=1}^{M}$
\REQUIRE Encoders $f_{\theta}$ (CNN), $g_{\phi}$ (MLP)
\REQUIRE Projection Heads $h_{\theta}$, $h_{\phi}$
\REQUIRE Batch Size $N$, Temperature $\tau$
\STATE Initialize parameters $\theta, \phi$ randomly
\FOR{\textbf{each} minibatch $\mathcal{B} = \{ (x_{\alpha}^{(k)}, x_{stat}^{(k)}) \}_{k=1}^{N}$ from $\mathcal{D}$}
    \STATE \textit{// Forward Pass}
    \STATE $v_{cnn} \leftarrow f_{\theta}(x_{\alpha})$ \COMMENT{Deep Semantic Features}
    \STATE $v_{mlp} \leftarrow g_{\phi}(x_{stat})$ \COMMENT{Deep Statistical Features}
    \STATE $z_{pay} \leftarrow h_{\theta}(v_{cnn})$, \quad $z_{stat} \leftarrow h_{\phi}(v_{mlp})$ \COMMENT{Project to latent space}
    \STATE
    \STATE \textit{// Compute Contrastive Loss}
    \STATE Compute similarity matrix $S = z_{pay} \cdot z_{stat}^T / (\tau \|z_{pay}\| \|z_{stat}\|)$
    \STATE $\mathcal{L}_{batch} \leftarrow \text{NT-Xent}(S)$ \COMMENT{Maximize agreement}
    \STATE
    \STATE \textit{// Update Weights}
    \STATE Update $\theta, \phi$ using Adam Optimizer to minimize $\mathcal{L}_{batch}$
\ENDFOR
\STATE \textbf{Return} Pre-trained CNN Encoder $f_{\theta^*}$
\end{algorithmic}
\end{algorithm}

\subsection{Phase II: Hybrid Late-Fusion \& Fine-Tuning}
\label{subsec:phase2_finetuning}

A critical architectural insight of LE-MVCL is the \textit{Hybrid Late-Fusion} strategy employed during the inference and fine-tuning stage. We empirically observed that projecting statistical features into a dense latent space (via the MLP) causes a loss of granular precision. Therefore, we propose a fusion mechanism that combines deep learned embeddings with raw statistical metrics.

\subsubsection{Fusion Mechanism}
Upon completion of pre-training, we discard the MLP encoder and the projection heads. We retain the pre-trained 1D-CNN encoder to extract robust semantic embeddings ($h_{cnn}$). The final feature vector for a flow is constructed by concatenating these embeddings with the normalized raw statistical components:

\begin{equation}
    \mathbf{v}_{final} = \text{Concat}\left( f_{\theta^*}(x_{\alpha}), \hat{x}_{\beta}, \hat{x}_{\gamma}, \hat{x}_{fft} \right)
\end{equation}

where $\hat{x}$ represents the Z-score normalized vectors of the flow dynamics, burst behavior, and spectral features. This hybrid approach ensures the model leverages the non-linear pattern recognition of the CNN (for obfuscation detection) while preserving the explicit resolution of metrics like Inter-Arrival Time (IAT) and burst duration, which are critical for distinguishing fine-grained VPN services (e.g., distinguishing Skype's VoIP patterns from Hangouts).

\subsubsection{Supervised Fine-Tuning}
The fused vector $\mathbf{v}_{final}$ is fed into an Extreme Gradient Boosting (XGBoost) classifier for the final decision. This ensemble method effectively handles the heterogeneity of the fused features. After extensive hyperparameter tuning, we selected the best-performing XGBoost configuration for the hybrid feature vector, as summarized in Table~\ref{tab:xgb_hyperparams}:

\begin{table}[htbp]
\centering
\caption{Best-performing XGBoost hyperparameters for the hybrid feature vector.}
\label{tab:xgb_hyperparams}
\begin{tabular}{l l}
\hline
\textbf{Hyperparameter} & \textbf{Value} \\
\hline
Estimators & 200 \\
Max Depth & 6 \\
Learning Rate & 0.05 \\
Subsample/Colsample & 0.8 \\
\hline
\end{tabular}
\end{table}

To handle class imbalance (e.g., rare VPN applications), we apply a \textit{Sample Weighting} strategy where weights are inversely proportional to class frequencies during the tree building process.
\section{Experimental Setup}
\label{sec:experimental_setup}

\subsection{Datasets and Evaluation Strategy}
\label{subsec:datasets_strategy}

To rigoroulsy assess the proposed LE-MVCL framework, we implemented a comprehensive evaluation protocol using the ISCX-VPN-2016 benchmark. Our strategy deviates from standard "flat" classification approaches by adopting a Cascaded Hierarchical Training Strategy \cite{parchekani2020rejecting,wang2025taonet}, designed to mirror real-world security operations where the primary challenge lies specifically within encrypted tunnels.

\subsubsection{The Cascaded Hierarchical Framework}
We treat the classification pipeline as two distinct, specialized modules:

\begin{enumerate}
    \item Module 1: The Encrypted Gatekeeper (Binary Task).
    The first module is trained on the complete dataset ($\mathcal{D}_{full}$), comprising both VPN and Non-VPN traffic. Its objective is to filter encrypted tunnels from standard background traffic. For this stage, we retain all available flows ($N \approx 12,555$) regardless of class size, ensuring the model learns to generalize across diverse traffic types.
    
    \item Module 2: The VPN Specialist (Fine-Grained Task).
    Once traffic is identified as VPN-encapsulated, it is passed to the second module. Crucially, this module is trained \textit{exclusively} on the VPN subset ($\mathcal{D}_{vpn}$). We isolate the VPN traffic to prevent the model from learning "shortcuts" based on unencrypted clear-text signatures (e.g., HTTP headers) found in Non-VPN traffic. This ensures that the reported accuracy reflects the model's true ability to characterize obfuscated tunnels.
\end{enumerate}

\subsubsection{Class Curation for Fine-Grained Tasks}
For the fine-grained experiments (Stage 2), we applied strict statistical curation to ensure validity.
\begin{itemize}
    \item Category Classification: We retained all 6 high-level categories (\textit{VoIP, Streaming, P2P, File Transfer, Email, Chat}). Even the smallest category (Email, 133 samples) provided sufficient data for stratified validation.
    \item Application Classification: To avoid misleading results driven by statistically insignificant classes, we applied a hard selection threshold: only applications with $>100$ flows were retained. Consequently, we focused our evaluation on 6 dominant applications: \textit{VPN\_Skype, VPN\_BitTorrent, VPN\_Hangout, VPN\_Facebook, VPN\_YouTube,} and \textit{VPN\_Email}. Minor classes such as \textit{VPN\_Netflix} (47 samples) or \textit{VPN\_SFTP} (11 samples) were excluded from the application-level experiment to prevent overfitting to outliers.
\end{itemize}

\subsubsection{Secondary Benchmark: VNAT (Robustness)}
To assess cross-domain generalization, we utilized the VNAT dataset solely for the Binary Task (Module 1). We employed a Zero-Shot Transfer strategy, training on ISCX and testing on VNAT without fine-tuning, to evaluate whether the learned distinction between "Encrypted" and "Standard" traffic holds across different network environments.

\subsubsection{Label Scarcity Protocol}
To validate our core contribution of Label Efficiency, we implemented a strict "Label Scarcity" training protocol. 
\begin{itemize}
    \item Phase I (Pre-training): The model is pre-trained on 100\% of the available training data \textit{without labels}, using the asymmetric contrastive objective.
    \item Phase II (Fine-tuning): The pre-trained encoder is frozen, and the fusion layer is fine-tuned using random stratified subsets of the labeled data, specifically at 5\%, 10\%, 20\%, and 30\% ratios.
\end{itemize}
This protocol allows us to directly measure the marginal gain provided by the pre-training phase compared to supervised baselines trained on the same limited subsets.

\subsubsection{Evaluation Metrics}
Given the inherent class imbalance in network traffic (e.g., P2P traffic often outweighs Email), we prioritize the Macro-Averaged F1-Score as our primary performance metric, alongside Precision and Recall.

\begin{equation}
    \text{F1}_{macro} = \frac{1}{C} \sum_{i=1}^{C} \frac{2 \cdot \text{Precision}_i \cdot \text{Recall}_i}{\text{Precision}_i + \text{Recall}_i}
\end{equation}

This ensures that each class contributes equally to the final score, preventing majority classes from dominating the reported performance and providing a fair assessment under severe label scarcity.

\subsection{Comparative Baselines}
\label{subsec:baselines}

To validate the three core pillars of our contribution—Label Efficiency, Asymmetric Alignment, and Direct Fusion—we benchmark LE-MVCL against three distinct categories of baselines. This comparative analysis is designed to isolate specific architectural decisions and prove their individual efficacy.

\subsubsection{Group A: Supervised Deep Learning (Label Efficiency Check)}
These baselines represent the current state-of-the-art in fully supervised learning. They are trained from scratch using the categorical cross-entropy loss.
\begin{itemize}
    \item 1D-CNN (End-to-End) \cite{wang2017end}: 
    The exact same convolutional architecture used in our LE-MVCL payload branch, but initialized with random weights and trained directly on the specific labeled subset (e.g., 10\%).
    \item Deep Packet \cite{lotfollahi2020deep}: 
    A standard Stacked Autoencoder (SAE) combined with a CNN, representing the established benchmark for encrypted traffic classification.
    \item \textit{Hypothesis Tested:} By comparing LE-MVCL to the standard \textit{1D-CNN}, we isolate the marginal gain provided specifically by the pre-training phase. A large gap here validates our solution to the label scarcity problem.
\end{itemize}

\subsubsection{Group B: Symmetric Contrastive Learning (The "Asymmetric" Check)}
To justify our novel "Asymmetric" strategy (aligning Payload $\leftrightarrow$ Stats), we compare it against a standard "Symmetric" Contrastive Learning approach typical of Computer Vision.
\begin{itemize}
    \item SimCLR-Net (Symmetric SSL) \cite{chen2020simple}: 
    A uni-modal framework where the contrastive objective aligns two augmented views of the \textit{Payload only}.
    \item \textit{Augmentation Strategy:} We employ "Byte Masking" (randomly zeroing out 10\% of the payload bytes) to generate the positive pairs.
    \item \textit{Hypothesis Tested:} This baseline tests the "Semantic Guide" theory. If LE-MVCL outperforms SimCLR-Net, it proves that using robust Flow Statistics as the positive anchor is superior to using artificially augmented payload views.
\end{itemize}

\subsubsection{Group C: Fusion Architectures (The "Direct Fusion" Check)}
Finally, we evaluate the impact of our fusion strategy on fine-grained classification.
\begin{itemize}
    \item Deep Fusion (Latent Concatenation) \cite{zeng2019deep}: 
    A variation of our model where the statistical features are passed through a non-linear MLP (Dense-64 $\rightarrow$ ReLU) before concatenation. The fusion occurs in the abstract latent space ($Z_{pay} \oplus Z_{stat}$).
    \item LE-MVCL (Ours - Direct Fusion): 
    Our proposed approach where \textit{raw, normalized} statistical features are concatenated directly with the CNN embeddings.
    \item \textit{Hypothesis Tested:} We hypothesize that Deep Fusion causes "manifold distortion," compressing vital statistical outliers. Superior performance by LE-MVCL would validate that preserving the granular precision of raw statistics is critical for distinguishing isotocol VPN applications.
\end{itemize}

\begin{table*}[htbp]
\caption{Summary of Comparative Baselines and Evaluation Goals}
\label{tab:baselines_summary}
\centering
\scriptsize
\setlength{\tabcolsep}{3pt}
\resizebox{\textwidth}{!}{%
\begin{tabular}{p{0.20\textwidth} p{0.20\textwidth} p{0.20\textwidth} p{0.32\textwidth}}
\hline
Baseline & Input Data & Learning Paradigm & Evaluation Goal \\
\hline 1D-CNN & Payload ($\alpha$) & Supervised (DL) & Check Label Efficiency \\
SimCLR-Net & Payload Only & Symmetric CL & Check "Semantic Guide" Theory \\
Deep Fusion & Payload + Stats & Late Fusion (MLP) & Check "Direct Fusion" Theory \\
LE-MVCL & Payload + Stats & Asymmetric CL & Proposed Framework \\
\hline
\end{tabular}%
}
\end{table*}
\section{Results and Discussion}
\label{sec:results}

\subsection{Evaluation of Label Efficiency}
\label{subsec:res_label_efficiency}

\begin{figure}[t]
    \centering
    \includegraphics[width=\columnwidth]{img/rank.png}
    \caption{Global Feature Importance. The top 20 features reveal a hybrid dependency, where Gamma (Burst) features provide the highest information gain, supported by learned Alpha (Payload) embeddings.}
    \label{fig:rank}
\end{figure}

\begin{figure}[t]
    \centering
    \includegraphics[width=\columnwidth]{img/featurecorrelationtop3.png}
    \caption{Linear Correlation Matrix. Beta features exhibit strong linear relationships, whereas the high-importance Gamma features are absent, highlighting the non-linear nature of burst dynamics.}
    \label{fig:correlation}
\end{figure}

\begin{figure*}[t]
    \centering
    \includegraphics[width=\textwidth]{img/violin.png}
    \caption{Multi-View Feature Distributions. Different views resolve different ambiguities: Beta isolates bulk P2P transfer, Gamma separates high-volume Streaming, and FFT uniquely identifies the periodic signature of VoIP.}
    \label{fig:violin}
\end{figure*}

\begin{figure}[t]
    \centering
    \includegraphics[width=\columnwidth]{img/spectral.png}
    \caption{Spectral Signature Profile. A parallel coordinates plot demonstrating the "Spectral Crossover": VoIP is characterized by high Periodicity (Max) and low Entropy, whereas P2P/Streaming exhibit the inverse.}
    \label{fig:spectral}
\end{figure}

\begin{figure}[t]
    \centering
    \includegraphics[width=\columnwidth]{img/spider.png}
    \caption{Holistic Traffic Fingerprints. Radar charts illustrate the unique multi-view geometry of each class, showing how LE-MVCL leverages orthogonal features to create separable signatures.}
    \label{fig:spider}
\end{figure}

\begin{figure*}[t]
    \centering
    \includegraphics[width=\textwidth]{img/tsne.png}
    \caption{Manifold Projection (t-SNE). Left: The baseline using only raw statistical features shows high confusion between interactive classes (VoIP/Chat). Right: LE-MVCL's Hybrid Fusion creates distinct, compact clusters, proving that the semantic "Alpha" embeddings successfully disentangle classes that are statistically similar.}
    \label{fig:tsne}
\end{figure*}

The primary hypothesis driving this research is that asymmetric contrastive pre-training can effectively mitigate the dependency on large-scale labeled datasets in encrypted traffic analysis. To validate this, we conducted a rigorous comparative evaluation against state-of-the-art supervised baselines, measuring performance degradation as label availability was restricted to 5\%, 10\%, 20\%, and 30\% of the training corpus.

To ensure a fair comparison of \textit{label efficiency}, we reproduced two established architectures: the 1D-CNN proposed by Wang et al. \cite{wang2017end} and Deep Packet by Lotfollahi et al. \cite{lotfollahi2020deep}. It is important to note the distinction in data access: supervised baselines are trained \textit{exclusively} on the labeled subset (e.g., 5\%) as they lack mechanisms to process unlabeled data. In contrast, LE-MVCL utilizes the remaining unlabeled portion of the dataset during the pre-training phase. This setup mirrors the real-world operational constraint where unlabeled traffic is abundant (free), but ground-truth labels are scarce (expensive). Thus, the comparison measures the \textit{value of leveraging unlabeled data} rather than pure sample efficiency.

\subsubsection{Performance Trajectory of LE-MVCL}
We first analyze the performance of the proposed LE-MVCL framework across the three hierarchical tasks: Binary Detection (Encrypted Gatekeeper), Category Classification, and Fine-Grained Application Identification. As shown in Table \ref{tab:lemvcl_trajectory}, LE-MVCL demonstrates robust scalability.

In the Binary task, the model achieves a strong F1-score of 89.8\% with only 5\% of the data, rapidly converging to 94.9\% at 30\% availability. This confirms the effectiveness of the pre-trained encoder in distinguishing fundamental encrypted tunnel structures from background traffic. For the more challenging fine-grained Application task, LE-MVCL maintains an F1-score of 75.2\% under extreme scarcity, which improves significantly to 91.3\% as more labels become available. This trajectory indicates that while the "cold start" performance is competitive, the model is highly responsive to additional supervision, effectively refining its decision boundaries for complex services like BitTorrent and Skype.

\begin{table}[htbp]
\caption{Performance Trajectory of LE-MVCL (Macro F1-Score)}
\label{tab:lemvcl_trajectory}
\centering
\begin{tabular}{l c c c c}
\hline
\textbf{Task} & \textbf{5\%} & \textbf{10\%} & \textbf{20\%} & \textbf{30\%} \\
\hline
Binary Detection & 89.8\% & 92.2\% & 94.5\% & 94.9\% \\
VPN Category & 70.7\% & 78.7\% & 84.8\% & 87.1\% \\
VPN Application & 75.2\% & 84.6\% & 90.0\% & 91.3\% \\
\hline
\end{tabular}
\end{table}

\subsubsection{Comparative Analysis with Baselines}
To benchmark these results, we compared LE-MVCL against four baselines: a lightweight statistical model (Stats-Only), two supervised deep learning models (1D-CNN, Deep Packet), and a symmetric contrastive model (SimCLR). Table \ref{tab:baseline_comparison} summarizes the performance gap ($\Delta$) between each baseline and LE-MVCL at the critical 5\% and 30\% thresholds.

\begin{table*}[htbp]
\caption{Comprehensive Baseline Comparison (Macro F1-Score). \\ Models marked with ($^\dagger$) utilize unlabeled data for pre-training (Semi-Supervised), while others rely solely on the labeled subset (Supervised).}
\label{tab:baseline_comprehensive}
\centering
\resizebox{\textwidth}{!}{%
\begin{tabular}{l | c c c c | c c c c | c c c c}
\hline
\multirow{2}{*}{\textbf{Model}} & \multicolumn{4}{c|}{\textbf{Binary Detection}} & \multicolumn{4}{c|}{\textbf{VPN Category}} & \multicolumn{4}{c}{\textbf{VPN Application}} \\
 & \textbf{5\%} & \textbf{10\%} & \textbf{20\%} & \textbf{30\%} & \textbf{5\%} & \textbf{10\%} & \textbf{20\%} & \textbf{30\%} & \textbf{5\%} & \textbf{10\%} & \textbf{20\%} & \textbf{30\%} \\
\hline
Stats-Only (XGB) & \textbf{91.0} & 92.0 & 93.0 & 94.0 & 72.0 & \textbf{79.0} & \textbf{87.0} & \textbf{89.0} & 74.0 & 83.0 & 87.0 & 88.0 \\
1D-CNN \cite{wang2017end} & 83.5 & 86.0 & 89.3 & 89.4 & 66.3 & 74.2 & 79.3 & 83.6 & 76.1 & 80.3 & 86.8 & 86.5 \\
Deep Packet \cite{lotfollahi2020deep} & 80.6 & 85.0 & 88.6 & 89.2 & 67.8 & 75.8 & 78.6 & 83.2 & 75.4 & 83.4 & 86.5 & 87.4 \\
Deep Fusion \cite{zeng2019deep} & 85.0 & 88.0 & 90.0 & 91.0 & \textbf{75.0} & 78.0 & 83.0 & 84.0 & \textbf{84.0} & 85.0 & 89.0 & 91.0 \\
\hline
SimCLR (Symmetric) $^\dagger$ & 88.4 & 91.9 & 93.1 & 93.5 & 65.2 & 78.1 & 82.2 & 85.8 & 70.4 & \textbf{85.5} & 88.2 & 89.5 \\
\textbf{LE-MVCL (Ours)} $^\dagger$ & 89.8 & \textbf{92.2} & \textbf{94.5} & \textbf{94.9} & 70.7 & 78.7 & 84.8 & 87.1 & 75.2 & 84.6 & \textbf{90.0} & \textbf{91.3} \\
\hline
\end{tabular}%
}
\end{table*}

Superiority in Binary Detection:
LE-MVCL consistently outperforms the deep learning baselines in the binary task. At 5\% labels, it exceeds the Supervised 1D-CNN by 6.3\% and Deep Packet by 9.2\%. While the Stats-Only model shows a slight initial advantage (+1.2\%), LE-MVCL overtakes it at 30\% (+0.9\%), proving that the learned payload representations provide scalability that pure statistics lack.

Scalability in Fine-Grained Tasks:
In the application identification task, LE-MVCL demonstrates superior long-term performance. At 30\% label density, it achieves the highest F1-score of 91.3\%, surpassing Stats-Only (88.0\%) and SimCLR (89.5\%). Although the Deep Fusion baseline shows strong initial performance in Application identification at 5\%, LE-MVCL closes the gap at 30\% (91.3\% vs 91.0\%), while consistently outperforming it in the Binary domain (89.8\% vs 85.0\% at 5\% labels), highlighting the robustness of our Direct Fusion strategy.

\subsection{Hierarchical Classification Performance}
\label{subsec:res_hierarchical}

We further evaluate the system's reliability by examining the class-wise performance of the "VPN Specialist" module. Table \ref{tab:hierarchical_breakdown} details the F1-scores for the top 6 VPN applications.

\begin{table}[htbp]
\caption{Stage 2: VPN Specialist Performance (Top-6 Apps, Macro F1)}
\label{tab:hierarchical_breakdown}
\centering
\scriptsize
\setlength{\tabcolsep}{3pt}
\resizebox{\columnwidth}{!}{%
\begin{tabular}{p{0.55\columnwidth} c c}
\hline
\textbf{Application} & \textbf{5\% Labels (F1)} & \textbf{30\% Labels (F1)} \\
\hline
VPN\_BitTorrent & 0.84 & 0.97 \\
VPN\_Skype & 0.86 & 0.97 \\
VPN\_Email & 0.79 & 0.92 \\
VPN\_YouTube & 0.76 & 0.95 \\
VPN\_Hangout & 0.73 & 0.87 \\
VPN\_Facebook & 0.54 & 0.80 \\
\hline
\textbf{Macro Average} & \textbf{0.75} & \textbf{0.91} \\
\hline
\end{tabular}%
}
\end{table}

The results show exceptional recovery for minority and complex classes. \textit{VPN\_BitTorrent}, often difficult due to its P2P nature, reaches a near-perfect 0.97 F1 at 30\%. \textit{VPN\_Skype} similarly hits 0.97. The most challenging class, \textit{VPN\_Facebook} (likely due to its multimodal mix of video and chat), sees a massive improvement from 0.54 to 0.80 as supervision increases, confirming that the model successfully disentangles it from similar Google services (YouTube/Hangout) when given sufficient examples.

\subsubsection{Cross-Domain Robustness (VNAT)}
Finally, to validate generalization, we tested the Binary Gatekeeper on the VNAT dataset (Table \ref{tab:binary_vnat}) using a Transfer Learning approach. The model achieved a sensitive Macro F1 of 98.33\% with only 5\% labeled data, demonstrating immediate adaptability. Performance remained stable across all label densities (98.24\% at 10\%, 97.99\% at 20\%, and 97.70\% at 30\%), proving that the learned features are robust and transferable across different network environments without extensive retraining.

\begin{table}[htbp]
\caption{Stage 1: Binary Gatekeeper Performance (VNAT Dataset)}
\label{tab:binary_vnat}
\centering
\scriptsize
\setlength{\tabcolsep}{3pt}
\resizebox{\columnwidth}{!}{%
\begin{tabular}{p{0.38\columnwidth} c c c c}
\hline
\textbf{Metric} & \textbf{5\% Labels} & \textbf{10\% Labels} & \textbf{20\% Labels} & \textbf{30\% Labels} \\
\hline
\textbf{Macro F1} & 98.33\% & 98.24\% & 97.99\% & 97.70\% \\
\hline
\end{tabular}%
}
\end{table}

\subsection{Multi-View Feature Analysis}
\label{subsec:feature_analysis}

To validate the LE-MVCL hypothesis—that encrypted traffic requires complementary feature views—we conducted a comprehensive analysis of feature importance, correlation, and spectral behavior.

\subsubsection{Global Feature Importance}
We first analyzed the global contribution of each feature view using XGBoost Information Gain. As shown in Fig. \ref{fig:rank}, the top 20 features are not dominated by a single view but are a hybrid mix of Gamma (Burst), Alpha (Payload), and Beta (Flow Stats). Notably, the single most critical feature is \texttt{gamma\_s2c\_burst\_pkts\_std} (Burst Packet Standard Deviation). This indicates that the variation in burst size acts as a primary "fingerprint" for distinguishing complex applications like Streaming and P2P.

\subsubsection{The Linearity-Utility Gap}
A critical finding emerges when comparing Feature Importance against the Linear Correlation Matrix (Fig. \ref{fig:correlation}). While Beta features (e.g., S2C Size Median) show strong linear correlations with classes like VoIP and Streaming, the top-performing Gamma features are notably absent from the top correlations. This reveals a \textit{"Linearity-Utility Gap,"} justifying the use of non-linear classifiers like XGBoost, as the most valuable burst dynamics serve as distinct thresholds rather than linear trends.

\subsubsection{Distributional Separability}
To visualize the decision boundaries created by these features, we examined their distributions in Fig. \ref{fig:violin}. The Beta View effectively isolates P2P traffic (bimodal size distribution), while the Gamma View distinguishes high-volume Streaming. The FFT View acts as a specialized "VoIP Detector," showing sharp vertical spikes (periodicity) unique to fixed-rate codecs.

\subsubsection{Spectral Analysis and Holistic Fingerprints}
Since time-domain statistics often overlap, we utilized the FFT view to analyze the frequency domain (Fig. \ref{fig:spectral}). We observed a "Spectral Crossover" where VoIP traffic exhibits distinct periodicity (Spectral Max) and low entropy, contrasting with the high-entropy "white noise" of bulk P2P. Finally, Fig. \ref{fig:spider} synthesizes these views into distinct class signatures, confirming that no two classes share the same geometric shape across all three views.

\subsubsection{Manifold Visualization}
To conclusively demonstrate the efficacy of the proposed fusion strategy, we projected the high-dimensional feature space into 2D using t-SNE (Fig. \ref{fig:tsne}). The baseline manifold (Left) exhibits significant overlap between VoIP and Chat. In contrast, the LE-MVCL manifold (Right) shows distinct, compact clusters with wide margins of separability, confirming that fusing semantic embeddings with granular statistics successfully disentangles statistically similar classes.




\section{Conclusion}
The rapid proliferation of encrypted tunneling protocols has rendered traditional deep packet inspection obsolete, creating a critical dependency on Deep Learning frameworks that are inherently data-hungry. In this paper, we presented LE-MVCL, a novel framework designed to break this dependency through \textit{Asymmetric Multi-View Contrastive Learning}.

Our work addresses the fundamental "duality" of network traffic by treating raw payload bytes and global flow statistics not as separate inputs, but as complementary views that guide each other. By aligning a deep 1D-CNN encoder with a robust statistical MLP during pre-training, we demonstrated that the statistical view acts as a "Semantic Guide," forcing the neural network to learn representations grounded in stable behavioral patterns rather than overfitting to volatile byte sequences. This approach yielded a decisive "Cold Start" advantage: LE-MVCL achieved a binary detection F1-score of 89.8\% using only 5\% of the labeled data, a threshold where standard supervised models effectively collapsed.

Furthermore, we empirically validated the superiority of Direct Statistical Fusion over standard Deep Fusion. Our results confirmed that projecting engineering metrics (like spectral entropy or burst duration) through deep non-linear layers results in "Statistical Compression," smoothing out the granular anomalies necessary for fine-grained classification. By preserving these features in their raw form, LE-MVCL successfully distinguished isotocol VPN applications—such as distinguishing \textit{BitTorrent} from \textit{Skype}—with a Macro F1-score of 91.3\%, proving that hybrid architectures can outperform pure end-to-end deep learning in domain-specific tasks.

While LE-MVCL sets a new benchmark for label efficiency, several avenues remain for future exploration:
\begin{itemize}
    \item Generalization to QUIC and ECH: As protocols evolve towards fully encrypted handshakes (e.g., TLS 1.3 Encrypted Client Hello and QUIC), the "Alpha" payload view may become less discriminative. Future work will investigate adapting the contrastive objective to focus more heavily on the \textit{temporal} structure of packet arrivals (e.g., using Flow Transformers) to compensate for the loss of header visibility.
    \item Open-World Classification: Our current framework assumes a closed-world setting with a fixed set of VPN applications. Extending LE-MVCL to an \textit{Open-World} scenario—where the model must detect and cluster "Unknown" encrypted flows without retraining—would significantly enhance its utility for zero-day anomaly detection in Security Operations Centers.
    \item Lightweight Edge Deployment: Although efficient in terms of labels, the dual-encoder architecture requires moderate computational resources. We plan to explore knowledge distillation techniques to compress the pre-trained LE-MVCL model into a lightweight student network suitable for deployment on resource-constrained edge routers or IoT gateways.
\end{itemize}

\section*{Acknowledgments}
This should be a simple paragraph before the References to thank those individuals and institutions who have supported your work on this article.







%{\appendices
%\section*{Proof of the First Zonklar Equation}
%Appendix one text goes here.
% You can choose not to have a title for an appendix if you want by leaving the argument blank
%\section*{Proof of the Second Zonklar Equation}
%Appendix two text goes here.}



% References
\bibliographystyle{IEEEtran}
\bibliography{references}


\newpage

\section{Biography Section}
If you have an EPS/PDF photo (graphicx package needed), extra braces are
 needed around the contents of the optional argument to biography to prevent
 the LaTeX parser from getting confused when it sees the complicated
 $\backslash${\tt{includegraphics}} command within an optional argument. (You can create
 your own custom macro containing the $\backslash${\tt{includegraphics}} command to make things
 simpler here.)
 
\vspace{11pt}

\bf{If you include a photo:}\vspace{-33pt}
\begin{IEEEbiography}[{\includegraphics[width=1in,height=1.25in,clip,keepaspectratio]{fig1}}]{Michael Shell}
Use $\backslash${\tt{begin\{IEEEbiography\}}} and then for the 1st argument use $\backslash${\tt{includegraphics}} to declare and link the author photo.
Use the author name as the 3rd argument followed by the biography text.
\end{IEEEbiography}

\vspace{11pt}

\bf{If you will not include a photo:}\vspace{-33pt}
\begin{IEEEbiographynophoto}{John Doe}
Use $\backslash${\tt{begin\{IEEEbiographynophoto\}}} and the author name as the argument followed by the biography text.
\end{IEEEbiographynophoto}




\vfill

\end{document}

